\documentclass[11pt]{article}
\usepackage{assets/tex/notes/eric-notes-style}

\begin{document}

\notetitle{MATH 327: Classical Mechanics Notes}{Newtonian, Lagrangian, and Hamiltonian Mechanics}

\topic{Newtonian Mechanics}

\entry{Newton's Law}
\[
F(t,x(t),\dot x(t))=\ddot x(t).
\]

\entry{Picard-Lindelof Thm}
\[
\dot x=G(t,x),\quad x(0)=x_0.
\]
\[
\text{If }G\text{ is differentiable, then }\forall x_0,V_0,\ \exists T>0
\text{ and a unique solution }x(t)
\]
\[
\text{for }-T<t<T,\ x(0)=x_0,\ \dot x(0)=v_0.
\]

\entry{Method using Taylor series}
\[
\text{For equation }G,\text{ we may write}
\]
\[
x(t)=\sum_{j=0}^{\infty}\frac{d^jx}{dt^j}(0)\frac{t^j}{j!}.
\]

\entry{Method from Euler-Rom scheme}
\[
\text{Let }x(t+h)=x(t)+\dot x(t)h,\qquad
x(t+2h)=x(t+h)+\dot x(t+h)h.
\]

\entry{Method from Picard iteration}
\[
\text{Let }x_0(t)=x(0),\qquad
x_{j+1}(t)=x_0(0)+\int_0^t G(s,x_j(s))\,ds.
\]

\entry{Conservation of Energy}
\[
F=-\nabla V(x).
\]
\[
\text{Given }F(x),\quad
1)\ \text{calculate }V(x)=-\int F(x)\,dx,
\qquad
E=\frac12\dot x^2+V(x).
\]
\[
2)\ \dot x=\sqrt{2(E-V(x))}.
\]
\[
3)\ \text{integrate to solve: }x(t)=x^{-1}(t+c).
\]

\entry{Equilibrium Points}
\[
\text{An equilibrium point }(x_0,0)\in\mathbb R^2\text{ satisfies }F(x_0)=0
\Longleftrightarrow V'(x_0)=0.
\]
\[
\text{It is stable if }V''(x_0)>0,\qquad
\text{unstable if }V''(x_0)<0.
\]

\entry{Phase portraits}
\[
\text{Given a potential }V(x),\text{ we can plot the phase portrait by}
\]
\[
\text{finding all equilibrium points and determining their stability.}
\]

\entry{Special case of Poincare Lemma}
\[
\nabla\times F=0,\ \text{i.e. } \partial_iF_j=\partial_jF_i,\quad
\forall i,j
\Longleftrightarrow
\exists\text{ smooth }V\in C^1\text{ s.t. }F=-\nabla V.
\]
\[
\text{If this holds, }F\text{ is conservative, and }
E=\frac12\|\dot x\|^2+V(x)\text{ is conserved.}
\]

\entry{Bound and unbound motion}
\[
\text{A trajectory }x:\mathbb R\to\mathbb R^n
\text{ is bound if }\exists R>0
\text{ s.t. }\|x(t)\|<R\ \forall t.
\]
\[
\text{Otherwise it is unbound.}
\]

\entry{Bound motion proposition}
\[
\text{Suppose }F:\mathbb R^n\to\mathbb R^n\text{ conservative, with potential }V.
\]
\[
\text{If }\exists R>0\text{ s.t. }\|x\|\le R\text{ for all }\|x\|>R,
\text{ then if }E<0,\text{ the motion is bound.}
\]

\newpage

\topic{Lagrangian Mechanics}

\entry{Lagrangian}
\[
\text{The Lagrangian }L=T-V,
\]
\[
\text{where }T\text{ is the kinetic energy of the system and }V
\text{ is the potential energy.}
\]

\entry{Stationary action}
\[
\gamma\text{ is a stationary point of the action}
\Longleftrightarrow
\text{Euler-Lagrange equation holds:}
\]
\[
L_x(t,\gamma,\dot\gamma)=\frac{d}{dt}L_v(t,\gamma,\dot\gamma).
\]

\entry{Noether's Thm}
\[
\text{Let }\gamma_T:x\mapsto x+w(x)+O(T^2)
\text{ be a continuous one parameter family,}
\]
\[
\text{where }w(x)=\left.\frac{d}{dT}\gamma_T(x)\right|_{T=0}
\quad\text{s.t.}\quad
L(Tx,D\gamma_T(x),v)=L(x,v).
\]
\[
\text{Then }p\cdot w=\frac{\partial L}{\partial v}\cdot w
\text{ is conserved.}
\]

\entry{Conservation of Hamiltonian}
\[
\text{If }L\text{ does not depend on time, then the Hamiltonian }H
\text{ is conserved.}
\]

\newpage

\topic{Hamiltonian Mechanics}

\entry{Legendre transform}
\[
\text{Given Lagrangian }L,\quad p_i=\frac{\partial L}{\partial \dot q_i},
\qquad \dot q_i=x_i\text{ where map }x(q,p).
\]
\[
\text{Also require }L_{ij}\ne0\Rightarrow p_i=0.
\]

\entry{Hamiltonian}
\[
\text{The Hamiltonian }H(q,p)=\sum_i p_i\dot q_i-L(q,\dot q).
\]

\entry{Hamilton's equations}
\[
\text{If }x(t)\text{ solves E-L, then }q(x(t)),p(x(t))
\text{ satisfy}
\qquad
\dot q_i=\frac{\partial H}{\partial p_i},\quad
\dot p_i=-\frac{\partial H}{\partial q_i}.
\]

\entry{Hamilton flow}
\[
\text{Hamilton }H\text{ generates a flow in phase space }\mathbb R^{2n},
\]
\[
X_H=\begin{pmatrix}\partial H/\partial p\\ -\partial H/\partial q\end{pmatrix},
\]
\[
\text{where }(q,p)\text{ evolves along an integral curve of the vector field.}
\]

\entry{Liouville's Thm}
\[
\text{Let }\phi_t:\mathbb R^{2n}\to\mathbb R^{2n}\text{ be the Hamiltonian flow.}
\]
\[
\int_{\phi_t(U)}1\,dV=\int_U1\,dV,\qquad
\text{i.e. }\phi\text{ is volume preserving.}
\]

\entry{Application in Stat Mech}
\[
\rho:\mathbb R^{2n}\to\mathbb R
\text{ is a density if }(i)\ \rho\ge0\text{ everywhere,}\quad
(ii)\ \int_{\mathbb R^{2n}}\rho(q,p)\,dq\,dp=1.
\]
\[
\rho\circ\phi_t\text{ is a density.}
\]
\[
\text{The flow generated by a vector field }V
\text{ is volume preserving }\Longleftrightarrow \nabla\cdot V=0.
\]

\entry{Poisson brackets}
\[
\text{Let }a(q,p)\in C^\infty(\mathbb R^n\times\mathbb R^n)
\text{ be an observable.}
\]
\[
\{q_i,H\}=\dot q_i,\qquad
\{p_i,H\}=\dot p_i,\qquad
\{f,g\}=-\{g,f\}.
\]
\[
\{f,g\}h+g\{f,h\}+\{h,f\}g=0.
\]
\[
\text{Can: if }\{H,H\}=0,\ \{f,H\}=0,\text{ then }\{f+g,H\}=0.
\]

\topic{Formula page}
\begin{itemize}
\item Tips paper
\item Spherical coordinates
\item Symmetry group
\end{itemize}

\end{document}
